\section{Machine Learning for Edge Scheduling}

More recent studies have explored using lightweight predictive models -- decision trees, gradient
boosted ensembles, and compact recurrent networks -- to forecast short-term environmental
variability and adjust sampling intervals accordingly. These approaches generally report energy
savings in the range of 15--30\% over fixed-rate baselines, though comparability across studies is
limited by inconsistent evaluation datasets and metrics.

Table~\ref{tab:related-work} summarises representative approaches from the reviewed literature.

\begin{table}[htbp]
\centering
\caption{Summary of related scheduling approaches}
\label{tab:related-work}
\begin{tabularx}{\textwidth}{@{}p{2.6cm}p{3.0cm}p{3.6cm}X@{}}
\toprule
\rowcolor{ummaroon!85}
\textcolor{white}{\textbf{Approach}} & \textcolor{white}{\textbf{Adaptivity}} &
\textcolor{white}{\textbf{Evaluation Data}} & \textcolor{white}{\textbf{Reported Limitation}} \\
\midrule
Fixed-rate duty cycling & None & Simulated traces & Ignores informational value of readings \\
\addlinespace
Threshold-based adaptive & Rule-based & Small pilot deployments & Manual per-site threshold tuning \\
\addlinespace
Regression-based forecasting & Statistical & Single-building datasets & Limited cross-building generalisation \\
\addlinespace
LSTM-based scheduling & Learned, sequential & Synthetic + short pilots & High inference cost on constrained gateways \\
\bottomrule
\end{tabularx}
\end{table}
